% Generado por docs/articulo/tabla_incluidos.py; no editar a mano.
\multicolumn{4}{@{}l}{\textbf{G1. Vías respiratorias altas, boca y ojos (19)}} \\
2012 & Infectious Diseases Society of America (IDSA)\newline{\footnotesize IDSA 2012: faringitis estreptocócica} & Internacional & \texttt{\scriptsize 10.1093/cid/cis629} \\
2012 & Korean Otologic Society\newline{\footnotesize Corea 2012: otitis media en niños} & Internacional & \texttt{\scriptsize 10.3346/jkms.2012.27.8.835} \\
2013 & American Academy of Pediatrics\newline{\footnotesize AAP 2013: otitis media aguda} & Internacional & \texttt{\scriptsize 10.1542/peds.2012-3488} \\
2013 & American Academy of Pediatrics\newline{\footnotesize AAP 2013: sinusitis bacteriana aguda de 1 a 18 años} & Internacional & \texttt{\scriptsize 10.1542/peds.2013-1071} \\
2019 & CODEPEH (España)\newline{\footnotesize CODEPEH 2019: otitis media con efusión} & Internacional & \texttt{\scriptsize 10.1016/j.otorri.2017.07.004} \\
2019 & Sociedad Chilena de Infectología (Comité de Infecciones Pediátricas)\newline{\footnotesize Chile 2019: faringoamigdalitis estreptocócica en pediatría} & Latinoamérica & \texttt{\scriptsize 10.4067/s0716-10182019000100069} \\
2019 & Sociedad Chilena de Infectología (Comité de Infecciones Pediátricas)\newline{\footnotesize Chile 2019: otitis media aguda en pediatría} & Latinoamérica & \texttt{\scriptsize 10.4067/s0716-10182019000400497} \\
2019 & Sociedad Chilena de Infectología (Comité de Infecciones Pediátricas)\newline{\footnotesize Chile 2019: rinosinusitis bacteriana aguda en pediatría} & Latinoamérica & \texttt{\scriptsize 10.4067/s0716-10182019000100078} \\
2019 & Società Italiana di Pediatria (SIP)\newline{\footnotesize SIP 2019: tratamiento de la otitis media aguda} & Internacional & \texttt{\scriptsize 10.1097/inf.0000000000002452} \\
2021 & DEGAM/DGHNO (Alemania)\newline{\footnotesize Alemania 2021: guía de dolor de garganta} & Internacional & \texttt{\scriptsize 10.3238/arztebl.m2021.0121} \\
2021 & Menzies School of Health Research y otros (Australia)\newline{\footnotesize Australia 2021: otitis media en niños aborígenes} & Internacional & \texttt{\scriptsize 10.5694/mja2.50953} \\
2022 & Hospital Infantil de México Federico Gómez y otros (México)\newline{\footnotesize México 2022: otitis media aguda en niños} & Latinoamérica & \texttt{\scriptsize 10.24875/bmhim.21000168} \\
2024 & Ministerio de Salud de Sudán\newline{\footnotesize Sudán 2024: fiebre reumática y cardiopatía reumática} & Internacional & \texttt{\scriptsize 10.3389/fcvm.2024.1403131} \\
2024 & SIPPS-SIP-SITIP-FIMP-SIAIP-SIMRI-FIMMG (Italia)\newline{\footnotesize Consenso italiano de faringitis 2024} & Internacional & \texttt{\scriptsize 10.1186/s13052-024-01789-5} \\
2025 & ABORL-CCF y otras (Brasil)\newline{\footnotesize Brasil 2024: rinosinusitis, evidencia y experiencia} & Latinoamérica & \texttt{\scriptsize 10.1016/j.bjorl.2025.101595} \\
2025 & Asociación Colombiana de Otorrinolaringología (ACORL)\newline{\footnotesize ACORL 2024: rinosinusitis en niños (Colombia)} & Colombia & \texttt{\scriptsize 10.37076/acorl.v52i4.808} \\
2025 & Consenso intersocietario italiano\newline{\footnotesize Consenso italiano de otitis media aguda 2025} & Internacional & \texttt{\scriptsize 10.1186/s13052-025-01894-z} \\
2025 & European Academy of Paediatric Dentistry (EAPD)\newline{\footnotesize EAPD 2025: uso de antibióticos en odontopediatría} & Internacional & \texttt{\scriptsize 10.1007/s40368-025-01068-9} \\
2025 & SIPPS-SIP-SITIP-FIMP-SIAIP-SIMRI-SIM-FIMMG (Italia)\newline{\footnotesize Consenso italiano de sinusitis 2025} & Internacional & \texttt{\scriptsize 10.1186/s13052-025-01868-1} \\
\addlinespace
\multicolumn{4}{@{}l}{\textbf{G2. Vías respiratorias bajas (15)}} \\
2011 & Pediatric Infectious Diseases Society e Infectious Diseases Society of America\newline{\footnotesize Guía PIDS/IDSA de neumonía adquirida en la comunidad en mayores de 3 meses} & Internacional & \texttt{\scriptsize 10.1093/cid/cir531} \\
2014 & Canadian Paediatric Society\newline{\footnotesize Bronquiolitis: recomendaciones de diagnóstico, monitorización y manejo (CPS)} & Internacional & \texttt{\scriptsize 10.1093/pch/19.9.485} \\
2017 & Saudi Pediatric Infectious Diseases Society\newline{\footnotesize Manejo de la neumonía adquirida en la comunidad en lactantes y niños (SPIDS)} & Internacional & \texttt{\scriptsize 10.1016/j.ijpam.2017.12.002} \\
2019 & Comité de Antimicrobianos, Sociedad Chilena de Infectología\newline{\footnotesize Recomendaciones para el diagnóstico y tratamiento antimicrobiano de la NAC bacteriana en pediatría (Chile)} & Latinoamérica & \texttt{\scriptsize 10.4067/s0716-10182019000400505} \\
2020 & South African Thoracic Society\newline{\footnotesize Diagnóstico y manejo de la neumonía adquirida en la comunidad en niños: actualización sudafricana 2020} & Internacional & \texttt{\scriptsize 10.7196/ajtccm.2020.v26i3.104} \\
2021 & Sociedad Argentina de Pediatría (comités de Neumonología, Infectología, Medicina Interna y Pediatría Ambulatoria)\newline{\footnotesize Recomendaciones para el manejo de las infecciones respiratorias agudas bajas en menores de 2 años, actualización 2021 (SAP)} & Latinoamérica & \texttt{\scriptsize 10.5546/aap.2021.s171} \\
2023 & Consenso intersocietario italiano\newline{\footnotesize Actualización 2022 de la guía italiana de bronquiolitis en lactantes} & Internacional & \texttt{\scriptsize 10.1186/s13052-022-01392-6} \\
2024 & Comisión de Pediatría, Sociedad Argentina de Infectología (SADI)\newline{\footnotesize Recomendaciones sobre neumonías complicadas en pediatría (Argentina)} & Latinoamérica & \texttt{\scriptsize 10.52226/revista.v32i116.345} \\
2024 & Grupo de expertos (China)\newline{\footnotesize Consenso de expertos sobre neumonía por Mycoplasma pneumoniae resistente a macrólidos en niños} & Internacional & \texttt{\scriptsize 10.1007/s12519-024-00831-0} \\
2024 & Grupo de expertos (China)\newline{\footnotesize Consenso de expertos sobre tos ferina en niños: nuevos conceptos en diagnóstico y tratamiento} & Internacional & \texttt{\scriptsize 10.1007/s12519-024-00848-5} \\
2024 & Grupo de expertos (China)\newline{\footnotesize Consenso de expertos sobre diagnóstico, tratamiento y prevención de la infección por VSR en niños} & Internacional & \texttt{\scriptsize 10.1007/s12519-023-00777-9} \\
2024 & SIPPS, SIP, SITIP, FIMP, SIAIP, SIMRI, FIMMG y SIMG (Italia)\newline{\footnotesize Tratamiento de la NAC leve a moderada en niños sanos: consenso intersocietario italiano} & Internacional & \texttt{\scriptsize 10.1186/s13052-024-01786-8} \\
2025 & Grupo de la guía nacional (Duodecim, Finlandia)\newline{\footnotesize Resumen de la actualización 2023 de la guía finlandesa de infecciones respiratorias bajas en niños} & Internacional & \texttt{\scriptsize 10.1111/apa.17481} \\
2025 & PREDICT (Australia y Nueva Zelanda)\newline{\footnotesize Guía australasiana de bronquiolitis: actualización 2025} & Internacional & \texttt{\scriptsize 10.1111/jpc.70144} \\
2025 & Sociedad de Pediatría de la Asociación Médica China\newline{\footnotesize Guía basada en la evidencia de neumonía por Mycoplasma pneumoniae en niños (2023)} & Internacional & \texttt{\scriptsize 10.1002/ped4.12469} \\
\addlinespace
\multicolumn{4}{@{}l}{\textbf{G3. Urinaria, piel, fiebre y meningitis (19)}} \\
2010 & University of Kentucky, Estados Unidos; publicada en Ann Saudi Med\newline{\footnotesize Guía de hospitalización por pielonefritis (Kentucky)} & Internacional & \texttt{\scriptsize 10.4103/0256-4947.68549} \\
2014 & Canadian Paediatric Society, Canadá\newline{\footnotesize CPS: meningitis bacteriana en mayores de 1 mes} & Internacional & \texttt{\scriptsize 10.1093/pch/19.3.141} \\
2014 & Comité Nacional de Infectología, SAP, Argentina\newline{\footnotesize SAP 2014: infecciones de piel y partes blandas (parte 2)} & Latinoamérica & \texttt{\scriptsize 10.5546/aap.2014.183} \\
2014 & Comité Nacional de Infectología, SAP, Argentina\newline{\footnotesize SAP 2014: infecciones de piel y partes blandas (parte 1)} & Latinoamérica & \texttt{\scriptsize 10.5546/aap.2014.96} \\
2017 & Sociedad Argentina de Pediatría, Argentina\newline{\footnotesize SAP: antibióticos en fiebre sin foco <36 meses (resumen ejecutivo)} & Latinoamérica & \texttt{\scriptsize 10.5546/aap.2017.205} \\
2018 & guías de IU pediátrica, evaluación AGREE II\newline{\footnotesize Revisión crítica (AGREE II) de guías de IU pediátrica} & Internacional & \texttt{\scriptsize 10.5489/cuaj.4796} \\
2019 & Países Bajos\newline{\footnotesize Guía clínica de Dientamoeba fragilis} & Internacional & \texttt{\scriptsize 10.1017/s0031182018001385} \\
2020 & Sociedad Italiana de Nefrología Pediátrica, Italia\newline{\footnotesize SINePe 2020: primera IU febril en niños de 2 meses a 3 años} & Internacional & \texttt{\scriptsize 10.1111/apa.14988} \\
2021 & SPIDS, Arabia Saudita\newline{\footnotesize SPIDS (Arabia Saudita): IU comunitaria de 90 días a 14 años} & Internacional & \texttt{\scriptsize 10.1016/j.ijpam.2021.03.001} \\
2021 & consenso suizo, Suiza\newline{\footnotesize Consenso suizo 2021 de IU en niños} & Internacional & \texttt{\scriptsize 10.1007/s00431-020-03714-4} \\
2021 & guía sueca comparada con AAP, CPS, NICE, España, Italia, EAU\newline{\footnotesize Comparación de la guía sueca de IU con otras} & Internacional & \texttt{\scriptsize 10.1111/apa.15727} \\
2022 & Comité de Nefrología, SAP, Argentina\newline{\footnotesize SAP 2022: guía de IU (actualización)} & Latinoamérica & \texttt{\scriptsize 10.5546/aap.2022.s69} \\
2022 & Sociedad Chilena de Infectología, Chile\newline{\footnotesize SOCHINF 2022: recomendaciones de IU (parte 1)} & Latinoamérica & \texttt{\scriptsize 10.4067/s0716-10182022000200174} \\
2022 & guías chinas e inglesas, AGREE II\newline{\footnotesize Evaluación AGREE II de guías de IU en niños} & Internacional & \texttt{\scriptsize 10.1136/bmjopen-2021-057736} \\
2023 & GPC colombiana, revista Infectio\newline{\footnotesize ACIN 2023: GPC colombiana de infecciones de piel y tejidos blandos} & Colombia & \texttt{\scriptsize 10.22354/24223794.1131} \\
2023 & guías 2000-2022\newline{\footnotesize Revisión narrativa de guías de IU en niños} & Internacional & \texttt{\scriptsize 10.4103/ua.ua\_147\_22} \\
2025 & guías de IU complicada\newline{\footnotesize Revisión integrativa de guías de IU complicada} & Internacional & \texttt{\scriptsize 10.1590/s1678-9946202567007} \\
2025 & guías hospitalarias de Londres y NICE NG143, Reino Unido\newline{\footnotesize Revisión de guías de Londres sobre lactantes febriles} & Internacional & \texttt{\scriptsize 10.1136/bmjpo-2024-002970} \\
2026 & EAU/ESPU, Europa\newline{\footnotesize EAU/ESPU 2025: IU en niños (resumen)} & Internacional & \texttt{\scriptsize 10.1016/j.jpurol.2025.06.016} \\
\addlinespace
\multicolumn{4}{@{}l}{\textbf{G4. Diarrea (5)}} \\
2017 & Groupe de Pathologie Infectieuse Pédiatrique, Sociedad Francesa de Pediatría (Arch Pediatr)\newline{\footnotesize Tratamiento antimicrobiano de la GEA en niños} & Internacional & \texttt{\scriptsize 10.1016/s0929-693x(17)30515-8} \\
2018 & IDSA y SHEA (Clin Infect Dis)\newline{\footnotesize Guía IDSA/SHEA de C. difficile, actualización 2017} & Internacional & \texttt{\scriptsize 10.1093/cid/cix1085} \\
2018 & Williams y Berkley (Paediatr Int Child Health)\newline{\footnotesize Tratamiento de la disentería (shigelosis): revisión sistemática} & OMS & \texttt{\scriptsize 10.1080/20469047.2017.1409454} \\
2020 & South African Society of Clinical Microbiology y otras sociedades (S Afr J Infect Dis)\newline{\footnotesize Guía SASCM de C. difficile} & Internacional & \texttt{\scriptsize 10.4102/sajid.v35i1.219} \\
2022 & Sociedad Pediátrica Jordana (Int J Pediatr)\newline{\footnotesize Consenso nacional jordano de GEA pediátrica} & Internacional & \texttt{\scriptsize 10.1155/2022/4456232} \\
\addlinespace
